\begin{frame}{경계의 물리적 이동 = 기저율 오류}
  \begin{itemize}
    \item 분류기가 소수 클래스에 대해 \textbf{보수적}이 됨.
    \item 3.1절 \kb{기저율 오류}와 정확히 같은 이야기:
      ``증거(가능도)가 아무리 강해도 사전확률이 낮으면 사후확률은
      기대만큼 안 올라간다''가 \textbf{결정 경계가 물리적으로
      이동하는} 형태로 나타남.
    \item 실전 클래스 비율 9:1이면 GDA는 ``소수 클래스로 예측하기
      전에 \textbf{더 확실한 증거}를 요구''하며 자동으로 경계를 옮겨줌.
    \item 실습 노트북에서는 2차원으로 시각화 --- 기울이는 같은 직선이
      $\phi$ 값에 맞게 \textbf{병렬 이동}.
  \end{itemize}
\end{frame}
